%% kapa: exam=1 round=22 year=2011 period=2 form=A subject=appraisal_law subject_ko=부동산관계법규 seq=40 number=80
%% source: Q-Net 감정평가사 기출문제 (cst00309) artlSeq=5036790 file=extracted/1008/5036790/2011년 제22회 감정평가사 1차시험 기출문제/2교시.hwp
%% conversion: kordoc 4.24.51 -> tools/gen1.py (auto); review=proofread-sonnet; flags=none
\begin{question}{80}
부동산등기법령상 등기에 관한 설명으로 옳지 않은 것은?
\begin{choices}
  \item 표제부의 등기는 이른바 사실의 등기이다.
  \item 미등기의 부동산에 대하여 처음으로 행하여지는 소유자의 등기는 보존등기이다.
  \item 어떤 등기가 행하여진 후 등기된 사항에 후발적 변경이 있어서 이를 바로잡기 위한 등기는 경정등기이다.
  \item 원시적 또는 후발적인 원인으로 등기와 실체관계의 전부가 부합하지 않을 경우에 이를 시정하기 위하여 등기내용의 전부를 소멸시키는 등기는 말소등기이다.
  \item 행정구역 또는 그 명칭이 변경되었을 때에는 등기부에 적은 행정구역 또는 그 명칭은 변경된 것으로 본다.
\end{choices}
\end{question}
\answer{3}
